\section{Rangkuman}
\label{sec:rangkuman-bab8}

Pada bab ini telah dibahas:

\begin{itemize}
    \item \term{Mode pengalamatan} adalah cara prosesor menyatakan
          letak operand. Ia menjadi jembatan antara penulisan
          simbolik programmer dan byte \eng{ModR/M} yang dihasilkan
          assembler.

    \item Tujuh mode pengalamatan 8086 adalah \eng{immediate},
          \eng{register}, \eng{direct}, \eng{register indirect},
          \eng{indexed}, \eng{base-relative}, dan
          \eng{based-indexed}.

    \item Mode \eng{immediate} dan \eng{register} tidak mengakses
          memori sehingga paling cepat. Mode \eng{direct} memakai
          offset tetap, sedangkan empat mode sisanya membentuk alamat
          dari register penunjuk.

    \item \term{Alamat efektif} dihitung sebagai jumlah basis,
          indeks, dan perpindahan: \code{EA = base + index + disp}.
          Alamat fisik kemudian dibentuk dengan menambahkan basis
          segmen, yaitu \code{(segmen x 16) + EA}.

    \item Hanya \reg{BX}, \reg{BP}, \reg{SI}, dan \reg{DI} yang boleh
          menjadi penunjuk alamat. Kombinasi seperti \code{[AX]},
          \code{[BX+BP]}, \code{[SI+DI]}, dan \code{[BX+SI+DI]}
          ditolak karena perangkat keras tidak menyediakan
          \eng{encoding}-nya.

    \item Segmen \eng{default} bergantung pada register penunjuk:
          kelompok \reg{DS} untuk \reg{BX}, \reg{SI}, dan \reg{DI};
          kelompok \reg{SS} untuk \reg{BP}. \eng{Segment override
          prefix} seperti \code{ES:} mengubah segmen yang dipakai
          tanpa mengubah nilai offsetnya.

    \item Efisiensi menurun seiring bertambahnya komponen alamat.
          Mode \eng{register} memerlukan dua siklus, sedangkan mode
          \eng{based-indexed} memerlukan sekitar sembilan belas
          siklus untuk instruksi yang sama panjangnya.

    \item Mode \eng{indexed} dipakai untuk menelusuri elemen array,
          sedangkan mode \eng{base-relative} dipakai untuk membaca
          medan \eng{record} dan parameter prosedur melalui
          \code{[BP+disp]}.
\end{itemize}

\begin{checklistbox}
Tandai setiap butir yang sudah Anda kuasai.

\begin{itemize}
    \item[\checkbox] Saya dapat menyebutkan tujuh mode pengalamatan
          8086 beserta contoh penulisannya.
    \item[\checkbox] Saya dapat menentukan mode pengalamatan hanya
          dari penulisan operandnya.
    \item[\checkbox] Saya dapat menghitung alamat efektif dan alamat
          fisik suatu operand.
    \item[\checkbox] Saya dapat menjelaskan mengapa \code{[AX]} dan
          \code{[SI+DI]} tidak diperbolehkan.
    \item[\checkbox] Saya dapat menentukan segmen \eng{default} dan
          menuliskan \eng{segment override prefix} dengan benar.
    \item[\checkbox] Saya dapat memilih mode yang paling efisien untuk
          mengakses elemen array dan medan \eng{record}.
    \item[\checkbox] Saya berhasil menelusuri alamat efektif dengan
          \eng{debugger} dan memverifikasinya pada jendela memori.
\end{itemize}

\vspace{0.4em}
Bila masih ada butir yang belum bertanda, ulangi bagian yang
bersesuaian sebelum melanjutkan ke Bab~\ref{chap:translasi}.
\end{checklistbox}

\begin{infobox}
\textbf{Rujukan bab.} Materi pada bab ini diperkuat oleh
\cite{intel1979} untuk daftar \eng{encoding} dan waktu siklus,
\cite{kader2014} untuk pengelompokan mode pengalamatan,
\cite{evans2020} untuk aturan segmen \eng{default} dan
\eng{override}, serta \cite{uva2024} untuk latihan penelusuran alamat
efektif. Untuk pembahasan segmen secara lebih mendalam, lihat
Bab~\ref{chap:segmentasi}.
\end{infobox}
